\title{DeepSeekMath: Pushing the Limits of Mathematical Reasoning in Open Language Models}

\begin{abstract}
The inherent complexity and structured requirements of mathematical reasoning present a formidable obstacle for language models. In this work, we present DeepSeekMath~7B, which extends the pre-training of DeepSeek-Coder-Base-v1.5 7B using a corpus of 120B math-focused tokens gathered from Common Crawl, complemented with both natural language and programming code datasets. DeepSeekMath~7B achieves a notable 51.7\% accuracy on the challenging MATH benchmark at the competition level, all without external tools or ensembling strategies, reaching a performance close to that of Gemini-Ultra and GPT-4. When employing self-consistency over 64 samples, DeepSeekMath~7B obtains a score of 60.9\% on the MATH benchmark. Two main factors underpin DeepSeekMath's proficiency in mathematical reasoning: First, we utilize the vast resources of publicly available web data through a carefully designed data selection workflow. Second, we propose Group Relative Policy Optimization (GRPO), an adaptation of Proximal Policy Optimization (PPO), which not only bolsters the model's mathematical reasoning but also improves the memory efficiency of PPO during optimization.
\end{abstract}

\section{Introduction}

Large language models (LLMs) have dramatically transformed the landscape of mathematical reasoning within artificial intelligence, driving remarkable progress in benchmarks like quantitative reasoning \citep{MATH} and geometry reasoning \citep{trinh2024solving}. Beyond these advancements, LLMs have shown considerable promise as tools to aid humans in solving sophisticated mathematical challenges \citep{tao}. Yet, high-performing systems such as GPT-4 \citep{gpt4} and Gemini-Ultra \citep{gemini} remain closed to the public, and the gap between these proprietary models and available open-source alternatives remains substantial.

This paper presents DeepSeekMath, a domain-focused language model that demonstrates markedly superior mathematical reasoning abilities compared to existing open-source models, approaching the proficiency of GPT-4 in various academic evaluations. To support this advancement, we compile the DeepSeekMath Corpus—a large-scale, high-caliber pre-training dataset containing 120B math tokens. The corpus is sourced from Common Crawl (CC) via a classifier based on fastText \citep{joulin2016fasttext}. During the first phase, the classifier is built using positive samples from OpenWebMath \citep{openwebmath}, supplemented by a heterogeneous selection of web content as negative instances. The classifier is subsequently deployed to extract further positive samples from CC, which are then meticulously filtered through manual annotation. This enriched labeled set is used to retrain the classifier, resulting in enhanced detection accuracy. Our evaluation demonstrates that the resulting dataset maintains high quality: the DeepSeekMath-Base 7B model attains a score of 64.2\% on GSM8K \citep{gsm8k} and 36.2\% on the competition-grade MATH benchmark \citep{MATH}, outperforming Minerva 540B \citep{minerva}. Moreover, due to the multilingual nature of the DeepSeekMath Corpus, we also observe substantial gains on Chinese math assessment sets \citep{wei2023cmath,agieval}. We contend that our efforts in processing mathematical datasets offer a valuable starting point for further investigation, acknowledging substantial avenues for future advancement.

DeepSeekMath-Base is constructed by initializing with DeepSeek-Coder-Base-v1.5 7B \citep{deepseek-coder}, as preliminary findings indicate that leveraging a code-centric pre-trained model yields superior results over initializing from a general-purpose LLM. Additionally, we observe that the subsequent math-specific training not only bolsters the model's proficiency in mathematics but also enhances performance on broader benchmarks such as MMLU \citep{mmlu} and BBH \citep{bbh}, suggesting that mathematical instruction contributes to improved general reasoning skills.

Following the pre-training phase, DeepSeekMath-Base undergoes mathematical instruction fine-tuning utilizing datasets comprising chain-of-thought \citep{cot}, program-of-thought \citep{pot,pal}, and tool-augmented reasoning \citep{tora} examples. The resultant model, DeepSeekMath-Instruct 7B, surpasses all other 7B-scale models in its category and demonstrates performance on par with instruction-tuned models of 70B scale from open-source repositories.

Moreover, we present Group Relative Policy Optimization (GRPO), a novel variant of reinforcement learning (RL) derived from the Proximal Policy Optimization (PPO) algorithm \citep{schulman2017proximal}. GRPO eliminates the need for a critic by estimating the learning baseline using group-based scoring, resulting in significantly lower resource requirements during training. With only a portion of English instruction tuning data, GRPO achieves marked improvements over the already strong DeepSeekMath-Instruct, attaining notable gains on both in-domain (GSM8K: 82.9\% $\rightarrow$ 88.2\%, MATH: 46.8\% $\rightarrow$ 51.7\%) and out-of-domain math datasets (e.g., CMATH: 84.6\% $\rightarrow$ 88.8\%) throughout the reinforcement learning stage. We also develop a unified framework to analyze a spectrum of approaches, including Rejection Sampling Fine-Tuning (RFT) \citep{yuan2023scaling}, Direct Preference Optimization (DPO) \citep{dpo}, PPO, and GRPO. Within this unified lens, these methodologies are characterized as forms of direct or simplified RL. Comprehensive ablation studies are conducted, examining factors such as online versus offline training, outcome versus process supervision, and single-turn compared to iterative RL, to thoroughly probe the critical aspects underpinning this framework. Finally, we articulate the mechanisms by which our RL approaches enhance instruction-tuned model performance and outline prospective avenues for developing more effective RL strategies within this unified paradigm.

\subsection{Contributions}
This work introduces advancements in large-scale mathematical pre-training, as well as a systematic study of reinforcement learning approaches.

\textbf{Math Pre-Training at Scale}

\begin{itemize}[topsep=0pt]
    \item We demonstrate that Common Crawl, an openly available data source, encompasses rich content suitable for mathematical machine learning. Leveraging a rigorously engineered data selection process, we compile the DeepSeekMath Corpus—a curated dataset comprising 120B tokens extracted from web sources with an emphasis on mathematical material. This dataset significantly surpasses previous collections, being approximately sevenfold greater than the corpus employed by Minerva \citep{minerva} and nine times larger than the recently introduced OpenWebMath \citep{openwebmath}.

    \item Our foundational model, DeepSeekMath-Base 7B, reaches performance levels comparable to the Minerva 540B model \citep{minerva}. This result challenges the prevailing notion that scaling up parameters is the primary determinant of a model's mathematical reasoning abilities. Rather, our findings suggest that a well-trained, smaller model with high-quality data can excel in such tasks.

    \item We present insights gleaned from our mathematical training studies, revealing that pre-training on code before math-specific data enhances the model's mathematical problem-solving, both with and without external tool support. This outcome provides partial evidence towards answering the long-debated question: \textit{Does code training enhance reasoning capabilities?} According to our observations, it does, at least for tasks involving mathematical reasoning.

    \item While training on arXiv materials remains prevalent, especially in the mathematical domain, our experiments find it yields no significant benefit across any of the mathematical benchmarks explored in this study.
\end{itemize}

\textbf{Exploration and Analysis of Reinforcement Learning}

\begin{itemize}[topsep=0pt]
    \item We present Group Relative Policy Optimization (GRPO), a novel reinforcement learning method that is both resource-efficient and effective. Unlike Proximal Policy Optimization (PPO), GRPO eliminates the need for a critic network by computing the baseline using group-level reward estimates, which leads to a substantial reduction in the computational demands during training.
    \item Our findings show that GRPO substantially boosts the capabilities of our instruction-tuned model, DeepSeekMath-Instruct, using exclusively instruction-tuning datasets. Additionally, we find improvements in out-of-domain generalization throughout the reinforcement learning phase.
    \item We develop a coherent framework for interpreting a range of methods, including RFT, DPO, PPO, and GRPO. Our experimental investigation covers multiple settings—such as online versus offline optimization, outcome-oriented versus process-based supervision, and single-turn as opposed to iterative reinforcement learning—to rigorously examine the crucial components underpinning this framework.
    \item Drawing from this integrated framework, we analyze the mechanisms that contribute to the effectiveness of reinforcement learning and propose several prospective avenues to further enhance the training of LLMs.
\end{itemize}

\subsection{Summary of Evaluations and Metrics}

\begin{itemize}[topsep=0pt]
    \item \textbf{English and Chinese Mathematical Reasoning}: Our evaluation spans comprehensive English and Chinese mathematical benchmarks, covering problems from elementary to undergraduate levels. The English suite includes GSM8K \citep{gsm8k}, MATH \citep{MATH}, SAT \citep{llemma}, OCW Courses \citep{minerva}, and MMLU-STEM \citep{mmlu}; for Chinese, we incorporate MGSM-zh \citep{mgsm}, CMATH \citep{wei2023cmath}, Gaokao-MathCloze \citep{agieval}, and Gaokao-MathQA \citep{agieval}. We assess the ability of our models both to generate complete textual solutions without external tools and to solve problems leveraging Python.

    On the English evaluations, DeepSeekMath-Base performs on par with Minerva 540B \citep{minerva}, a proprietary system, and notably outperforms all existing open-source base models—such as Mistral 7B \citep{mistral} and Llemma-34B \citep{llemma}—regardless of their pre-training regimen, often by a large margin. Importantly, DeepSeekMath-Base excels on Chinese tasks, which we attribute to our inclusion of diverse, high-quality multilingual training data instead of exclusively English sources as seen in prior studies \citep{minerva,llemma}. Through both mathematical instruction tuning and reinforcement learning, our enhanced models DeepSeekMath-Instruct and DeepSeekMath-RL achieve robust results, including surpassing the 50\% accuracy threshold on the challenging MATH benchmark—a first for an open-source approach.
\end{itemize}

\item \textbf{Formal Mathematics}: We assess the abilities of DeepSeekMath-Base on the informal-to-formal theorem proving task described in \citep{dsp_proof}, utilizing the miniF2F benchmark \citep{minif2f} with Isabelle \citep{isabelle} as the proof assistant. Our evaluation indicates that DeepSeekMath-Base achieves notable results in few-shot autoformalization.

\item \textbf{Natural Language Understanding, Reasoning, and Code}: To comprehensively gauge the model’s capacities in general language understanding, reasoning, and programming, we test DeepSeekMath-Base on several benchmarks: Massive Multitask Language Understanding (MMLU) \citep{mmlu}, which covers 57 varied multiple-choice domains; BIG-Bench Hard (BBH) \citep{bbh}, comprised of 23 advanced tasks mainly involving multi-step reasoning; and HumanEval \citep{codex} as well as MBPP \citep{mbpp}, both established benchmarks for code language models. We find that pre-training with mathematical data enhances both the model’s language comprehension and reasoning abilities.
\end{itemize}

\section{Math Pre-Training}

\subsection{Data Collection and Decontamination} This section describes the methodology for assembling the DeepSeekMath Corpus from Common Crawl. Illustrated in Figure \ref{fig:data_collect}, our iterative pipeline details a systematic approach to curating a substantial collection of mathematical data, beginning with a compact but reliable seed dataset of math-related resources. It should be emphasized that the same procedure can be adapted for other fields, such as code.

To initiate, we select OpenWebMath \citep{openwebmath}, recognized for its high-quality mathematical web texts, as the seed dataset. Leveraging this dataset, we train a fastText model \citep{joulin2016fasttext} to retrieve web pages from Common Crawl that closely resemble the OpenWebMath style. Specifically, we extract 500,000 examples at random from the seed dataset as positive samples and supplement them with 500,000 unrelated web pages from Common Crawl as negative samples. Training is conducted using an open-source fastText library\footnote{\url{https://fasttext.cc}} with the following parameters: vector dimensionality of 256, a learning rate of 0.1, a word n-gram maximum length of 3, a minimum occurrence of 3 per word, and 3 training epochs. To manage the scale of Common Crawl, we apply URL-based deduplication as well as near-duplicate removal, resulting in a refined dataset of 40 billion HTML pages. The fastText model is then employed to surface mathematical content from this deduplicated collection. For quality assurance, we rank candidate pages based on their fastText-predicted scores and select only the highest-ranked subset. The quantity of data retained is further evaluated through pre-training trials on the top 40B, 80B, 120B, and 160B tokens. In our initial cycle, we opt to keep the top 40B tokens.

Following the initial data collection round, a significant portion of mathematical web pages remains uncaptured. This is primarily due to the limited diversity within the positive samples used to train the fastText model. To address this, we broaden the range of mathematical web sources, thereby augmenting the initial seed corpus and facilitating improved fastText optimization. Our methodology begins by partitioning the entire Common Crawl dataset into distinct domains, each defined by web pages that share an identical base URL. For every domain, we compute the fraction of pages collected during the first pass. If more than 10\% of a domain’s pages have been retrieved, we categorize that domain as math-focused (for instance, \url{mathoverflow.net}). 

Subsequent steps involve manually labeling URLs within these math-designated domains that are explicitly tied to mathematical material (such as \url{mathoverflow.net/questions}). We then expand the seed corpus by including any uncollected web pages linked to these URLs. This incremental expansion enhances the range of positive training examples, thereby enabling the fastText model to recognize and extract a larger volume of mathematical content in future iterations. Through four iterative rounds of data collection, our process results in the aggregation of 35.5 million mathematical web pages, amounting to a total of 120 billion tokens. By the fourth round, we observe that nearly 98\% of mathematical content had already been gathered in the preceding round, leading us to discontinue further data acquisition.

To mitigate the risk of benchmark data leakage, our filtering protocol mirrors the approach described in \cite{deepseek-coder}. Specifically, we excise any web page containing questions or answers present in established English-language mathematical benchmarks—such as GSM8K \citep{gsm8k} and MATH \citep{MATH}—as well as Chinese benchmarks like CMATH \citep{wei2023cmath} and AGIEval \citep{agieval}. The procedure entails removing any segment from the training corpus in which a 10-gram substring perfectly matches any substring from an evaluation benchmark. For evaluation materials shorter than 10 grams but at least 3 grams in length, we similarly exclude web pages based on exact string matches.

\subsection{Validating the Quality of the DeepSeekMath~Corpus}

To evaluate the utility of the DeepSeekMath~Corpus, we conduct a series of pre-training experiments benchmarking its performance against several contemporary math-centric corpora:

\begin{itemize}[topsep=0pt]
    \item \textbf{MathPile} \citep{mathpile}: A composite dataset of 8.9B tokens amassed from textbooks, Wikipedia, ProofWiki, CommonCrawl, StackExchange, and arXiv, with the latter representing the predominant source (over 85\%).
    \item \textbf{OpenWebMath} \citep{openwebmath}: Curated from CommonCrawl and filtered to contain 13.6B tokens of mathematical material.
    \item \textbf{Proof-Pile-2} \citep{llemma}: Comprised of OpenWebMath, AlgebraicStack (including 10.3B tokens of mathematical code), and arXiv documents (amounting to 28.0B tokens). Experiments with Proof-Pile-2 follow the 2:4:1 arXiv:Web:Code ratio as specified in \cite{llemma}.
\end{itemize}

\subsubsection{Training Setting}
\label{sec:quality-policy}
We fine-tune a general-purpose pre-trained language model with 1.3 billion parameters, utilizing the same architecture as DeepSeek LLMs \citep{deepseek-llm}, referred to as DeepSeek-LLM 1.3B in this work. The model undergoes separate training on each mathematical dataset for a total of 150 billion tokens. All training runs make use of the efficient HAI-LLM \citep{haillm} infrastructure. Consistent with the procedures outlined for DeepSeek LLMs, we adopt the AdamW optimizer \citep{adamW} configured with $\beta_1=0.9$, $\beta_2=0.95$, and $\mathrm{weight\_decay}=0.1$. The learning rate schedule is multi-staged: after a warmup phase of 2{,}000 steps where the learning rate increases to its maximum, it is reduced to 31.6\% of this peak after 80\% of the training duration, and then further to 10.0\% of the peak after completing 90\% of training. The maximum learning rate is set at 5.3e-4. The batch size is fixed at 4 million tokens, and the maximum sequence length per context window is set to 4{,}000 tokens.

\subsubsection{Evaluation Results}

\textbf{The DeepSeekMath~Corpus is distinguished by its quality, scale, and multilingual mathematical data coverage.}

\begin{itemize}[topsep=0pt]
    \item \textbf{High-quality}:
    The evaluation comprises downstream tasks from 8 different mathematical benchmarks using few-shot chain-of-thought prompting \cite{cot}. Table \ref{tab:corpora_comparison} presents results where the model trained on the DeepSeekMath~Corpus clearly outperforms its counterparts. Further, Figure \ref{fig:corpora_comparisons} demonstrates that training on the DeepSeekMath Corpus yields consistently superior outcomes compared to Proof-Pile-2 when measured at the 50B token mark (which corresponds to one complete pass over Proof-Pile-2), implying a higher average sample quality in DeepSeekMath Corpus.
    \item \textbf{Multilingual}:
    Data within the DeepSeekMath~Corpus is drawn from numerous languages, with English and Chinese being most prevalent. As demonstrated in Table \ref{tab:corpora_comparison}, models trained on this corpus achieve notable gains in mathematical reasoning across both English and Chinese, unlike existing predominantly English-only mathematical datasets, which offer minimal benefits or even a decrease in Chinese task performance.
    \item \textbf{Large-scale}:
    The scale of the DeepSeekMath~Corpus surpasses that of all currently available mathematical datasets by a considerable margin. According to Figure \ref{fig:corpora_comparisons}, DeepSeek-LLM 1.3B trained with DeepSeekMath~Corpus exhibits a more pronounced and sustained increase in learning performance. Conversely, the smaller baselines have been subjected to multiple training epochs, which results in a rapid performance plateau.
\end{itemize}

\subsection{Training and Evaluating DeepSeekMath-Base 7B}

This section details the development of DeepSeekMath-Base 7B, a foundational model characterized by robust reasoning skills with a particular emphasis on mathematics. We start with DeepSeek-Coder-Base-v1.5 7B \citep{deepseek-coder} as the initialization point, subsequently training it on a total of 500B tokens. The training data is distributed as follows: 56\% originates from the DeepSeekMath Corpus, 4\% from AlgebraicStack, 10\% from arXiv, 20\% from Github code repositories, while the remaining 10\% comprises natural language data from Common Crawl in both English and Chinese. The majority of the training protocol adheres to the setup outlined in Section \ref{sec:quality-policy}, with two exceptions: we utilize a peak learning rate of 4.2e-4 and a batch size consisting of 10M tokens.

To gauge the model’s mathematical proficiency, we perform a thorough evaluation emphasizing several facets: generation of fully self-contained mathematical solutions without external resources, solving problems with tool assistance, and executing formal theorem proofs. Additionally, the evaluation encompasses general capabilities of the base model, such as its competencies in natural language understanding, complex reasoning, and coding.

\paragraph{Mathematical Problem Solving with Step-by-Step Reasoning}

We assess DeepSeekMath-Base’s mathematical problem-solving abilities employing few-shot chain-of-thought prompting \citep{cot} across eight distinct benchmarks in both English and Chinese. These datasets include tasks requiring quantitative reasoning—such as GSM8K \citep{gsm8k}, MATH \citep{MATH}, and CMATH \citep{wei2023cmath}—as well as multiple-choice assessments like MMLU-STEM \citep{mmlu} and Gaokao-MathQA \citep{agieval}. Collectively, these benchmarks span a broad spectrum of mathematical topics, ranging from elementary to collegiate complexity levels.

Table \ref{tab:base_math_cot} illustrates that DeepSeekMath-Base 7B achieves the highest scores across all eight benchmarks compared to other open-source base models, including the broadly adopted Mistral 7B \citep{mistral} and the more recent Llemma 34B \citep{llemma}, which has been trained on mathematical datasets like Proof-Pile-2 \citep{llemma}. Particularly striking is its performance on the MATH dataset, representative of competition-level challenges, where DeepSeekMath-Base 7B exceeds all prior open-source base models by more than 10\% in absolute terms, and even outperforms Minerva 540B \citep{minerva}—a proprietary base model that is 77 times larger, built upon PaLM \citep{palm}, and additionally trained on mathematical literature.

\paragraph{Mathematical Problem Solving with Tool Use} We assess program-assisted mathematical reasoning on GSM8K and MATH by applying few-shot program-of-thought prompting techniques \citep{pot,pal}. In these tasks, models are instructed to solve problems by generating Python scripts, employing packages like \textit{math} and \textit{sympy} as needed for complex computations. The correctness of the solution is determined by executing the produced code and comparing its output with the expected answer. As indicated in Table \ref{tab:base_math_pot_proof}, DeepSeekMath-Base 7B establishes a new state-of-the-art among open models, surpassing Llemma 34B in this setting.

\paragraph{Formal Mathematics}

Automating formal proofs contributes to the correctness, reliability, and efficiency of mathematical proof construction, and has garnered substantial interest in recent times. To evaluate DeepSeekMath-Base 7B on this front, we consider the informal-to-formal proving challenge introduced in \citep{dsp_proof}. This involves synthesizing a formal proof from a natural language theorem statement, its formal translation, and a corresponding informal proof. Using the miniF2F benchmark \citep{minif2f}, which features problems at the level of mathematical Olympiads, we employ few-shot prompting to generate Isabelle-formalized proofs. Following the procedure in \cite{dsp_proof}, the model produces proof outlines which are then completed by the automated tool Sledgehammer \citep{sledgehammer}. Results, presented in Table \ref{tab:base_math_pot_proof}, demonstrate that DeepSeekMath-Base 7B excels in automated formalization of mathematical reasoning.

\paragraph{Natural Language Understanding, Reasoning, and Code}

We further benchmark model capabilities in natural language understanding on MMLU \citep{mmlu}, reasoning skills on BBH \citep{bbh}, and coding proficiency using HumanEval \citep{codex} and MBPP \citep{mbpp}. Table \ref{tab:understanding_reasoning_code} displays that DeepSeekMath-Base 7B achieves substantial improvements on MMLU and BBH relative to its predecessor, DeepSeek-Coder-Base-v1.5 \citep{deepseek-coder}, confirming the benefits of incorporating mathematics in model training for language understanding and reasoning tasks. With the introduction of code tokens during continued training, DeepSeekMath-Base 7B also preserves the robust coding performance observed in DeepSeek-Coder-Base-v1.5 on the HumanEval and MBPP tasks. Across all three domains—reasoning and coding—DeepSeekMath-Base 7B consistently and significantly surpasses the general-purpose Mistral 7B \citep{mistral}.

\section{Supervised Fine-Tuning}

\subsection{SFT Data Curation}

We assemble a dataset specifically for mathematical instruction-tuning, encompassing both English and Chinese questions that span a broad array of mathematical domains and difficulty levels. Each problem is paired with corresponding solutions presented in chain-of-thought (CoT) \citep{cot}, program-of-thought (PoT) \citep{pot,pal}, and tool-integrated reasoning formats \citep{tora}. This collection comprises a total of 776K training samples.

\begin{itemize}[topsep=0pt]
    \item \textbf{English mathematical datasets}: For English-language data, we annotate items from GSM8K and MATH with tool-assisted solutions, and incorporate a selection from MathInstruct \citep{MathInstruct} as well as the Lila-OOD \citep{lila} training set, where each question is solved using either CoT or PoT approaches. This English dataset spans various mathematical disciplines such as algebra, probability, number theory, calculus, and geometry.
    \item \textbf{Chinese mathematical datasets}: For the Chinese portion, we aggregate K-12 math problems distributed over 76 sub-categories (e.g., linear equations), providing solutions annotated both in CoT and tool-integrated formats.
\end{itemize}

\subsection{Training and Evaluating DeepSeekMath-Instruct 7B}

Here, we describe the development and assessment of DeepSeekMath-Instruct 7B, which has undergone supervised instruction-tuning on the DeepSeekMath-Base model. During fine-tuning, individual training samples are concatenated at random, up to a maximum input length of 4K tokens. The training process consists of 500 optimization steps, using a batch size of 256 and a fixed learning rate set at 5e-5.

The mathematical reasoning ability of models is assessed both in settings without tool support and with tool augmentation, across four benchmarks focused on quantitative reasoning in both English and Chinese. We compare the performance of our approach against prominent state-of-the-art models available at the time:

\begin{itemize}[topsep=0pt]
    \item \textbf{Closed-source models} include: (1) various iterations of the GPT family, among which GPT-4 \citep{gpt4} and GPT-4 Code Interpreter~\footnote{\url{https://openai.com/blog/chatgpt-plugins\#\#code-interpreter}} demonstrate top-tier performance, (2) Gemini Ultra and Pro \citep{gemini}, (3) Inflection-2 \citep{inflection-2}, (4) Grok-1~\footnote{\url{https://x.ai/model-card}}, in addition to recently introduced models from Chinese tech companies such as (5) Baichuan-3~\footnote{\url{https://www.baichuan-ai.com}} and (6) the most recent GLM-4~\footnote{\url{https://open.bigmodel.cn/dev/api\#glm-4}}
    from the GLM family \citep{glm}.
\end{itemize}

The general-purpose nature of these models is underscored by their comprehensive alignment processes.
\item \textbf{Open-source models} encompass: broadly applicable models such as (1) DeepSeek-LLM-Chat 67B \citep{deepseek-llm}, (2) Qwen 72B \citep{qwen}, (3) SeaLLM-v2 7B \citep{seallm}, and (4) ChatGLM3 6B \citep{chatglm3}, in addition to several models explicitly designed to enhance mathematical performance, including (5) InternLM2-Math 20B~\footnote{\url{https://github.com/InternLM/InternLM-Math}}, which is based on InternLM2 with supplementary mathematical training followed by instruction tuning; (6) Math-Shepherd-Mistral 7B, leveraging PPO training \citep{schulman2017proximal} on top of Mistral 7B \citep{mistral} using a process-oriented reward model; (7) the WizardMath family \citep{wizardmath}, which elevates the mathematical proficiency of Mistral 7B and Llama-2 70B \citep{llama2} using evolve-instruct—a technique involving AI-generated instructions—and PPO, with primary training data drawn from GSM8K and MATH; (8) MetaMath 70B \citep{metamath}, created by fine-tuning Llama-2 70B with expanded GSM8K and MATH datasets; (9) ToRA 34B \cite{tora}, derived from CodeLlama 34B and optimized for tool-augmented mathematical reasoning; and (10) MAmmoTH 70B \citep{MathInstruct}, which involves instruction tuning Llama-2 70B on MathInstruct.

As depicted in Table \ref{tab:sft_rl_math}, within the framework that restricts tool usage, DeepSeekMath-Instruct 7B exhibits notable proficiency in performing detailed step-by-step reasoning. On the challenging MATH competition benchmark, our model outperforms every available open-source competitor as well as most commercial systems (including Inflection-2 and Gemini Pro) by a minimum margin of 9 percentage points, regardless of whether those models are much larger in scale (such as Qwen 72B) or specifically enhanced with mathematical reinforcement learning methods (such as WizardMath-v1.1 7B). DeepSeekMath-Instruct demonstrates performance on par with Chinese commercial models GLM-4 and Baichuan-3 for the MATH dataset, but does not yet reach the performance of GPT-4 and Gemini Ultra.

In scenarios where models are permitted to combine natural language reasoning with code-based tool usage during problem solving, DeepSeekMath-Instruct 7B attains close to 60\% accuracy on MATH, exceeding all prior open-source solutions. Across additional benchmarks, its results are comparable to DeepSeek-LLM-Chat 67B—the previously leading open-source model which is an order of magnitude larger.

\section{Reinforcement Learning}

\subsection{Group Relative Policy Optimization}
Reinforcement learning (RL) techniques are known to provide substantial gains in LLMs' mathematical reasoning abilities following the Supervised Fine-Tuning (SFT) phase \citep{wang2023math,wizardmath}. Here, we present Group Relative Policy Optimization (GRPO), our RL algorithm tailored for both effectiveness and efficiency.

\subsubsection{Transitioning from PPO to GRPO}
Proximal Policy Optimization (PPO) \citep{schulman2017proximal} stands out as a widely adopted actor-critic RL method, frequently utilized during RL fine-tuning of LLMs \citep{ouyang2022training}. PPO improves model parameters by maximizing the following surrogate objective function:

\begin{equation}
\footnotesize
    \mathcal{J}_{PPO}(\theta) = \mathbb{E}{[q \sim P(Q), o \sim \pi_{\theta_{old}}(O|q)]} \frac{1}{|o|} \sum_{t=1}^{|o|} \min \left[ \frac{\pi_\theta(o_{t} | q, o_{<t})}{\pi_{\theta_{old}}(o_{t} | q, o_{<t})} A_{t}, \text{clip} \left( \frac{\pi_\theta(o_{t} | q, o_{<t})}{\pi_{\theta_{old}}(o_{t} | q, o_{<t})}, 1 - \varepsilon, 1 + \varepsilon \right)  A_{

In this context, $\pi_{\theta}$ and $\pi_{\theta_{old}}$ denote the present and previous versions of the policy model, while $q$ and $o$ represent questions and the corresponding outputs, respectively, sampled from the question set and generated by the previous policy $\pi_{\theta_{old}}$. The parameter $\varepsilon$ serves as a hyper-parameter for clipping, playing a crucial role in the Proximal Policy Optimization (PPO) framework by helping to stabilize the training process. The advantage, $A_t$, is determined through Generalized Advantage Estimation (GAE) \citep{gae}, utilizing reward signals $\{r_{\ge t}\}$ and predictions from a learned value function $V_{\psi}$. Consequently, PPO necessitates training a value function model concurrently with the policy model. Furthermore, to curb excessive reward model optimization, it is standard practice to incorporate a per-token KL divergence penalty against a reference model into the reward at each generation step, as recommended by \citep{ouyang2022training}:

\begin{equation}
    r_{t} = r_\varphi(q, o_{\le t}) - \beta \log\frac{\pi_{\theta}(o_{t}|q, o_{<t})}{\pi_{ref}(o_{t}|q, o_{<t})},
\label{eq:PPO-reward}
\end{equation}

Here, $r_\varphi$ is the reward model, $\pi_{ref}$ designates the reference model (often initialized as the SFT model), and $\beta$ is the scaling factor for the KL divergence penalty.

However, because the value function employed by PPO is generally as large as the policy model, this introduces significant computational and memory costs. Moreover, during the reinforcement learning phase, the value function serves primarily as a baseline to reduce variance in the advantage estimates. In the realm of large language models (LLMs), reward models often only evaluate the final token, rather than providing rewards for every token, making it challenging to train a value function that yields accurate predictions at each timestep. To alleviate this, as illustrated in Figure \ref{fig:grpo}, we introduce Group Relative Policy Optimization (GRPO). GRPO removes the requirement for a separate value function approximator, as required in PPO. Instead, it constructs a baseline using the mean reward of multiple outputs generated for a single question. Specifically, for each question $q$, GRPO draws a group of outputs $\{o_1, o_2, \cdots, o_G\}$ from $\pi_{\theta_{old}}$, then updates the policy by maximizing the following criterion:

\begin{equation}
\footnotesize
\begin{split}
    \mathcal{J}_{GRPO}(\theta) &= \mathbb{E}{[q \sim P(Q), \{o_i\}_{i=1}^G \sim \pi_{\theta_{old}}(O|q)]}  \\
    & \frac{1}{G}\sum_{i=1}^G\frac{1}{|o_i|} \sum_{t=1}^{|o_i|} \left\{ \min \left[ \frac{\pi_\theta(o_{i,t} | q, o_{i,<t})}{\pi_{\theta_{old}}(o_{i,t} | q, o_{i,<t})} \hat{A}_{i,t}, \text{clip} \left( \frac{\pi_\theta(o_{i,t} | q, o_{i,<t})}{\pi_{\theta_{old}}(o_{i,t} | q, o_{i,<t})}, 1 - \varepsilon, 1 + \varepsilon \right)  \hat{A}_{i,t} \right] - \beta \mathbb{D}_{KL}\left[\pi_{\theta} || \pi_{ref}\right]\right\} ,
\end{split}
\label{eq:GRPO-obj}
\end{equation}

Here, both $\varepsilon$ and $\beta$ denote hyper-parameters, and $\hat{A}_{i,t}$ corresponds to the advantage, which is now defined by comparing rewards exclusively among the outputs within each group—a procedure further described in subsequent sections. This group-relative approach is particularly well-suited to the comparative design of reward models, since such models are usually trained with pairs of outputs for the same prompt. Moreover, GRPO directly penalizes the KL divergence between the policy and reference models in the loss term, rather than incorporating the KL into the reward, thereby simplifying the estimation of $\hat{A}_{i,t}$. Distinct from the KL penalty in (\ref{eq:PPO-reward}), we estimate the

\begin{algorithm}[t]
  \small
  \caption{Iterative Group Relative Policy Optimization}
  \textbf{Input} initial policy model $\pi_{\theta_{\text{init}}}$; reward models $r_\varphi$; task prompts $\mathcal{D}$; 
  hyperparameters $\varepsilon$, $\beta$, $\mu$
  \begin{algorithmic}[1]
    \State Initialize the policy model: $\pi_\theta \leftarrow \pi_{\theta_{\text{init}}}$
    \For{iteration = 1, \dots, I}
       \State Set the reference model $\pi_{ref} \leftarrow \pi_{\theta}$
      \For{step = 1, \dots, M}
      \State Draw a batch $\mathcal{D}_b$ from $\mathcal{D}$
      \State Set the old policy: $\pi_{\theta_{old}} \leftarrow \pi_{\theta}$ 
      \State For each question $q \in \mathcal{D}_b$, sample $G$ responses $\{o_i\}_{i=1}^G \sim \pi_{\theta_{old}} (\cdot \mid q)$
      \State Obtain reward values $\{r_i\}_{i=1}^{G}$ for each sampled output $o_i$ using $r_{\varphi}$ 
      \State For every token $t$ in $o_i$, calculate $\hat{A}_{i,t}$ using group-based relative advantage estimation.
      \For{GRPO iteration = 1, \dots, $\mu$}
        \State Adjust $\pi_\theta$ to increase the GRPO objective (refer to Equation \ref{eq:GC-GRPO})
      \EndFor
    \EndFor \State Refine the reward model $r_\varphi$ continually with a replay training mechanism. 
    \EndFor 
  \end{algorithmic}
  \textbf{Output} $\pi_\theta$
  \label{alg:iter-grpo}
\end{algorithm}

\subsubsection{Outcome Supervision RL with GRPO} In the formal setup, for any question $q$, a set of responses $\{o_1, o_2, \cdots, o_G\}$ is generated from the prior policy $\pi_{\theta_{old}}$. Each output is scored by a reward model, resulting in $G$ reward values $\mathbf{r}=\{r_1, r_2, \cdots, r_G\}$. These reward scores undergo normalization, where the group mean is subtracted and the result divided by the group’s standard deviation. Using outcome supervision, the final token of every output $o_i$ receives the normalized reward, which is then broadcast as the advantage $\hat{A}_{i, t}$ to every token in $o_i$; thus, $\hat{A}_{i, t} = \widetilde{r}_i = \frac{r_i- {\rm mean}(\mathbf{r})}{{\rm std}(\mathbf{r})}$. The policy is then updated by maximizing the objective in equation (\ref{eq:GRPO-obj}).

\subsubsection{Process Supervision RL with GRPO} While outcome supervision offers a terminal reward for each sequence, this sparse signal can be inadequate for guiding policies on sophisticated mathematical problems. Following the method outlined in \cite{wang2023math}, process supervision is adopted, delivering rewards at the completion of each intermediate reasoning segment. More precisely, given a question $q$ and a collection of $G$ outputs $\{o_1, o_2, \cdots, o_G\}$, a process-oriented reward model assesses every intermediate step, producing sets of rewards: $\mathbf{R} = \{ \{r_1^{index(1)},\cdots,r_1^{index(K_1)}\}, \cdots,  \{r_G^{index(1)},\cdots,r_G^{index(K_G)}\} \}$, where $index(j)$ refers to the last token of the $j$-th step, and $K_i$ is the number of steps in the $i$-th response. Each process-level reward is normalized similarly: $\widetilde{r}_i^{index(j)} = \frac{r_i^{index(j)} - {\rm mean}(\mathbf{R})}{{\rm std}(\mathbf{R})}}$.
The advantage for a token is then determined by summing the normalized rewards from all subsequent steps: $\hat{A}_{i, t} = \sum_{index(j) \ge t} \widetilde{r}_i^{index(j)}$.

\subsubsection{Iterative RL with GRPO} Over the course of reinforcement learning training, the previously trained reward model may no longer provide adequate guidance for the policy model's current capabilities. To address this, we further investigate iterative RL utilizing GRPO. As outlined in Algorithm \ref{alg:iter-grpo}, the iterative GRPO procedure involves creating fresh datasets for the reward model by drawing samples from the evolving policy model. The reward model is then continually updated through a replay strategy, wherein 10\% of prior data is incorporated into each update cycle. Subsequently, the policy model is designated as the new reference model and is iteratively optimized with feedback from the updated reward model.

\subsection{Training and Evaluating DeepSeekMath-RL}

Our reinforcement learning experiments are performed using DeepSeekMath-Instruct 7B as the policy foundation. The RL training set comprises approximately 144K chain-of-thought-style items, selected from the GSM8K and MATH components within the SFT dataset. To specifically evaluate RL's effects on benchmarks lacking direct training data, questions unrelated to GSM8K and MATH are omitted. The reward model training data is assembled according to the protocol described in \citep{wang2023math}. Initial reward model training begins with DeepSeekMath-Base 7B and employs a learning rate of 2e-5. For GRPO updates, we adopt a policy learning rate of 1e-6 and set the KL regularization coefficient at 0.04. Each input question receives $64$ sampled outputs. Both the maximum input length and batch size are fixed at 1024. The policy network is updated once per exploration cycle. DeepSeekMath-RL 7B's performance is assessed on the same suite of benchmarks as DeepSeekMath-Instruct 7B. For DeepSeekMath-RL 7B, tasks such as GSM8K and MATH—with chain-of-thought solutions—are categorized as in-domain, while remaining benchmarks are considered out-of-domain tasks.

Table \ref{tab:sft_rl_math} reports the results for open- and closed-source models employing chain-of-thought and tool-augmented reasoning strategies across English and Chinese tasks. Our analysis reveals the following: 1) DeepSeekMath-RL 7B achieves 88.2\% and 51.7\% accuracy on GSM8K and MATH, respectively, using chain-of-thought approaches. These outcomes not only exceed those of all other open-source models in the 7B–70B category but also outperform most closed-source alternatives. 2) Notably, DeepSeekMath-RL 7B is exclusively trained on chain-of-thought instruction data from GSM8K and MATH, with DeepSeekMath-Instruct 7B as its initial model. Even with such limited data, it consistently surpasses DeepSeekMath-Instruct 7B in every evaluation metric, underscoring the value of reinforcement learning.

\section{Discussion}
This section presents insights gleaned from our pre-training and reinforcement learning (RL) investigations.

\subsection{Insights from Pre-Training}

We begin by outlining observations from our pre-training process. Unless specified otherwise, the following accounts assume the training framework detailed in Section \ref{sec:quality-policy}. For reference, the DeepSeekMath Corpus mentioned here corresponds to an 89B-token version assembled during the second phase of data collection.

\subsubsection{Incorporating Code Improves Mathematical Reasoning}

An often-discussed but largely unproven assumption is that code training enhances a model’s reasoning capabilities. We provide empirical evidence supporting this notion within mathematical contexts: incorporating code training bolsters a model's proficiency in mathematical reasoning, whether or not external tools are used.

To investigate the influence of code training on mathematical reasoning, we conducted experiments under two paradigms: two-stage training and one-stage training. 

\textbf{Two-Stage Training}

\begin{itemize}[topsep=0pt]
    \item \textbf{Code Training for 400B Tokens $\rightarrow$ Math Training for 150B Tokens}: DeepSeek-LLM 1.3B undergoes pretraining on 400B code tokens, after which it is further trained on 150B math tokens.
    \item \textbf{General Training for 400B Tokens $\rightarrow$ Math Training for 150B Tokens}: To serve as a control, we replace the initial code tokens with 400B general tokens drawn from DeepSeek-AI’s extensive general corpus, allowing us to analyze whether code-centric training confers unique advantages for mathematical reasoning tasks over training with general content.
\end{itemize}

\textbf{One-Stage Training}

\begin{itemize}[topsep=0pt]
    \item \textbf{Math Training for 150B Tokens}: DeepSeek-LLM 1.3B is directly trained for 150B tokens solely on mathematics data.
    \item \textbf{Training on a mixture of 400B Code Tokens and 150B Math Tokens}: Since introducing math training after code pretraining can harm code proficiency, we examine whether a combined training regimen—incorporating both 400B code tokens and 150B math tokens within a single training phase—can simultaneously enhance mathematical reasoning and mitigate catastrophic forgetting.
\end{itemize}

\paragraph{Results}
Empirical findings are reported in Table \ref{tab:code-math} and Table \ref{tab:code-math-general-eval}, summarizing downstream task results for each training regime.

Training with code tokens demonstrably advances program-based mathematical reasoning under both multi-phase and mixed-phase conditions. Table \ref{tab:code-math} illustrates that, within a two-phase process, code pretraining alone leads to marked gains in Python-enabled problem-solving for datasets such as GSM8K and MATH, with subsequent math-specific training bringing additional progress. Remarkably, a single-phase approach blending code and math tokens curtails catastrophic forgetting associated with staged training, while also boosting both coding capacity (see Table \ref{tab:code-math-general-eval}) and performance in program-supported mathematical reasoning (Table \ref{tab:code-math}).

Furthermore, code-centric training provides advantages in mathematical reasoning tasks that do not employ external tools. During two-stage training, early exposure to code imparts moderate improvements and paves the way for more effective subsequent math learning, culminating in superior overall results. By contrast, a one-stage strategy that intermingles code and math tokens negatively impacts pure mathematical reasoning. We hypothesize that the 1.3B-parameter DeepSeek-LLM may be too small to concurrently learn from and internalize both code and math modalities at sufficient depth within a single stage.

\subsubsection{ArXiv Papers Appear Unhelpful for Advancing Mathematical Reasoning}

ArXiv papers are a prevalent source within the math pre-training datasets \citep{minerva,gpt-f,llemma,mathpile}. Yet, a thorough evaluation of their contribution to mathematical reasoning has remained largely unexplored. Somewhat surprisingly, our results suggest that arXiv papers offer little to no benefit in fostering mathematical reasoning abilities. We assess this across models of varying sizes, namely DeepSeek-LLM 1.3B and DeepSeek-Coder-Base-v1.5 7B \citep{deepseek-coder}, leveraging arXiv datasets subjected to distinct preprocessing workflows:

\begin{itemize}[topsep=0pt]
    \item \textbf{MathPile} \citep{mathpile}: a corpus containing 8.9B tokens, curated using cleaning and filtering heuristics, where scientific arXiv papers comprise over 85\% of the data;
    \item \textbf{ArXiv-RedPajama} \citep{redpajama}: a full collection of arXiv LaTeX source files, stripped of preambles, comments, macros, and references, totaling 28.0B tokens.
\end{itemize}

For our study, DeepSeek-LLM 1.3B was trained for 150B tokens, while DeepSeek-Coder-Base-v1.5 7B was trained for 40B tokens on each arXiv dataset. The evidence indicates that exposure to arXiv content does not enhance mathematical reasoning capabilities. Models trained solely on arXiv material demonstrated stagnant or even declining performance across a spectrum of math benchmarks of varied difficulty applied in our evaluation. The tested benchmarks encompass quantitative reasoning datasets such as GSM8K and MATH (Table \ref{tab:arxiv-cot}), multiple-choice tasks including MMLU-STEM (Table \ref{tab:arxiv-cot}), and formal mathematics challenges such as miniF2F (Table \ref{tab:arxiv_atp}).

Nonetheless, the generalizability of this finding is limited. Several aspects remain unexplored, notably:

\begin{itemize}[topsep=0pt]
    \item Potential advantages of arXiv data for targeted mathematical tasks not addressed here, such as the informalization of theorems—i.e., translating formal arguments or results into their informal articulations;
    \item The interaction effects of arXiv data when incorporated with heterogeneous datasets;
    \item Whether scaling to larger model architectures might reveal benefits of arXiv paper data.
\end{itemize}

These open questions suggest further investigations are warranted and are beyond the present study’s scope.

\subsection{Insights of Reinforcement Learning}
\subsubsection{Toward a Unified Paradigm} We introduce a unifying framework to examine a range of training strategies—including SFT, RFT, DPO, PPO, and GRPO—and perform empirical analyses to disentangle critical factors underlying this framework. In general, the gradient of a training method with respect to the parameter $\theta$ can be formulated as:

\begin{equation}
    \nabla_{\theta}\mathcal{J}_{\textcolor{red}{\mathcal{A}}}(\theta) = \mathbb{E}[\underbrace{(q,o) \sim \textcolor{red}{\mathcal{D}}}_{Data \ Source}]\left( \frac{1}{|o|} \sum_{t=1}^{|o|}  \underbrace{GC_{{\mathcal{A}}}(q, o, t, \textcolor{red}{\pi_{{rf}}})}_{Gradient \ Coefficient}  \nabla_{\theta}\log \pi_{\theta}(o_t | q, o_{<t})\right).
\label{eq:objective}
\end{equation}

This formulation encompasses three principal elements: 1) \textit{Data Source $\mathcal{D}$}, specifying the corpus from which training samples are drawn; 2) \textit{Reward Function $\pi_{{rf}}$}, furnishing the training signal through its evaluation of model behavior; and 3) \textit{Algorithm $\mathcal{A}$}, responsible for converting the training samples and associated reward signals into the gradient coefficient $GC$ that ultimately controls the strength and direction of updates to the model's parameters. The following summarizes notable techniques interpreted within this unifying schema:

\begin{itemize}
    \item  \textbf{Supervised Fine-tuning (SFT)}: SFT further adapts a pretrained model by training it on datasets manually curated by humans.
    \item \textbf{Rejection Sampling Fine-tuning (RFT)}: RFT additionally trains the SFT model using outputs sampled from the SFT model itself, filtered by whether their answers to given prompts are correct.
    \item \textbf{Direct Preference Optimization (DPO)}: DPO improves upon SFT by fine-tuning on expanded output sets created by the SFT model, leveraging a pairwise loss that incorporates relative preferences.
    \item \textbf{Online Rejection Sampling Fine-tuning (Online RFT)}: In contrast to standard RFT, Online RFT initializes from the SFT model but then continually updates the policy model by fine-tuning on outputs sampled from the most recent version of the policy itself.
    \item \textbf{PPO/GRPO}: PPO/GRPO commence with the SFT-trained model as the policy's initialization, and subsequently enhance it using outputs sampled from the evolving policy model during training.
\end{itemize}

The elements constituting these methods are outlined in Table \ref{tab:data_policy}. For a comprehensive derivation and deeper discussion, see Appendix \ref{app:analysis-rl}.

\paragraph{Discussion of Data Source} We categorize data sources into online and offline sampling. Online sampling refers to acquiring training data from trajectories generated by the current policy model during learning, whereas offline sampling relies on data produced by the pre-trained SFT model. The RFT and DPO algorithms utilize offline sampling, while Online RFT and GRPO are grounded in online sampling methodologies.

Referencing Figure \ref{fig:rl-analysis}, it becomes apparent that Online RFT delivers consistently superior results compared to RFT across two evaluation benchmarks. At the onset of training, performance between Online RFT and RFT is similar; however, as training progresses, Online RFT secures a notable performance edge, reflecting the benefits of online data collection. This result aligns with expectations: in early training phases, the actor and SFT models are nearly indistinguishable, meaning the sampled datasets are largely alike. As training continues, distinctions in samples drawn from the actor model become more pronounced, rendering real-time sampling increasingly advantageous.

\paragraph{Discussion of Gradient Coefficient} The reward information is incorporated into the gradient coefficient, which steers the parameter updates in the learning algorithm. In our experiments, we employ both `Rule' and `Model' based reward functions. The `Rule' reward assesses responses directly based on answer accuracy, while the `Model' reward leverages a trained reward model, itself trained on data labeled using rule-based judgments, to assign scores to responses. As detailed in Equations \ref{eq:GC-RFT} and \ref{eq:GC-GRPO}, GRPO and Online RFT diverge in a key aspect: GRPO dynamically modifies its gradient coefficient according to the score returned by the reward model. This mechanism enables more nuanced positive and negative reinforcement that corresponds to the reward magnitude. By contrast, Online RFT does not possess this adaptive capability; it fails to penalize incorrect answers, providing uniform positive reinforcement to all correct responses, regardless of score distinctions.

As illustrated in Figure \ref{fig:rl-analysis}, GRPO outperforms online RFT, underscoring the advantage of dynamically modulating the positive and negative gradient coefficients. Moreover, GRPO+PS yields better results than GRPO+OS, which underscores the value of leveraging more granular, step-sensitive gradient coefficients. We additionally investigate iterative RL, conducting two iterations in our experiments. According to Figure \ref{fig:iter-rl}, this iterative process leads to substantial performance gains, particularly noticeable after the first iteration.

\subsubsection{Why RL Works?} 
In this study, reinforcement learning is applied using a subset of the instruction tuning data, resulting in significant improvements over the baseline instruction-tuned model. To shed light on the underlying mechanisms of reinforcement learning's effectiveness, we assess both Pass@K and Maj@K accuracies for the Instruct and RL models across two evaluation datasets. Figure \ref{fig:maj-pass} reveals that while RL leads to enhanced Maj@K accuracy, Pass@K remains largely unaffected. This suggests that reinforcement learning primarily strengthens the robustness of the output distribution; that is, \textbf{the observed improvement appears to stem from an increased likelihood of generating the correct answer within the TopK responses, rather than a direct boost in intrinsic model competence.} In a similar vein, \citep{wang2023making} reported a \textbf{misalignment problem} in SFT models on reasoning tasks and demonstrated that preference alignment strategies could bolster the reasoning ability of such models \citep{yuan2023rrhf,song2023preference,wang2023making}.

\subsubsection{How to Achieve More Effective RL?}
\label{sec:effec_rl}
Our findings establish that RL is highly effective for mathematical reasoning challenges. We further introduce a unified framework to interpret various major training approaches, positioning them as forms or simplifications of RL. As encapsulated in Equation \ref{eq:objective}, this framework delineates three essential elements: Data Source, Algorithm, and Reward Function. We outline several avenues for advancing each of these components in the future.

\paragraph{Data Source} The data source serves as the foundation for any training method. Within RL, this specifically refers to unlabeled questions paired with outputs generated from the policy model. This work confines the data source to queries from the instruction tuning phase, employing simple nucleus sampling for output generation. We hypothesize that this limitation could explain why our RL pipeline primarily boosts Maj@K. Future work will involve expanding the RL process to include prompts beyond the original training distribution, together with \textbf{more sophisticated sampling (decoding) methods}, such as those utilizing tree-search strategies \citep{yao2023tree}. Moreover, advances in \textbf{efficient inference techniques} \citep{xia-etal-2023-speculative,leviathan2023fast,kwon2023efficient,xia2024unlocking}, which impact the exploration efficiency of policy models, are anticipated to play a crucial role.

\paragraph{Algorithms} Algorithms leverage both data and reward signals to calculate gradient coefficients, which in turn are used to adjust model parameters. According to Equation \ref{eq:objective}, existing approaches largely \textbf{TRUST} the feedback from the reward function, adjusting the conditional probability of specific tokens upward or downward in response. Nonetheless, there are inherent limitations to the reliability of reward signals, especially in scenarios involving highly intricate tasks. For instance, despite the rigorous annotation by trained professionals, the PRM800K datasets \citep{lightman2023let} are still estimated to exhibit around 20\% erroneous annotations\footnote{\url{https://github.com/openai/prm800k/issues/12\#issuecomment-1728491852}}. Motivated by this, we turn our attention to reinforcement learning algorithms that maintain robustness in the face of noisy rewards. We anticipate that \textbf{WEAK-TO-STRONG} alignment techniques \citep{burns2023weak} could fundamentally reshape the way learning algorithms operate.

\paragraph{Reward Function} The reward function acts as the origin of the training feedback signal. In reinforcement learning, this function is frequently implemented as a neural reward model. We identify three key areas for the advancement of reward models: 1) \textbf{Strategies to bolster the reward model’s capacity to generalize.} It is crucial for the reward model to perform reliably on out-of-distribution queries and novel generated responses; without this, reinforcement learning may simply reinforce the status quo of large language model distributions, rather than fostering substantive improvement; 2) \textbf{Mechanisms to represent and utilize the uncertainty within the reward model.} Proper modeling of uncertainty could facilitate the integration of imperfect reward models with weak-to-strong learning frameworks; 3) \textbf{Methods for the rapid construction of precise process reward models} that can deliver detailed and instructive feedback throughout complex reasoning chains \citep{lightman2023let,wang2023math}.

\section{Conclusion, Limitation, and Future Work}

We introduce DeepSeekMath, which surpasses all open-source alternatives on the competitive MATH benchmark and demonstrates results comparable to proprietary systems. DeepSeekMath~is initialized using DeepSeek-Coder-v1.5 7B, followed by continual training on 500B tokens, with 120B math-related tokens curated primarily from Common Crawl making up a large portion of the training set. Through comprehensive ablation experiments, we find that data gathered from web pages serves as a valuable resource for high-quality mathematical content, whereas arXiv contributions fall short of our initial expectations. Additionally, we propose Group Relative Policy Optimization (GRPO), a novel adaptation of Proximal Policy Optimization (PPO), which substantially enhances mathematical reasoning while requiring less memory. Experimental analysis demonstrates that GRPO maintains its benefits even as DeepSeekMath-Instruct 7B achieves strong benchmark performance. Furthermore, we establish a unified framework for interpreting various techniques and outline multiple promising research directions to foster more effective reinforcement learning.

Despite achieving strong results on benchmarks for quantitative reasoning, DeepSeekMath~displays relative shortcomings in areas such as geometry and theorem-proving compared to proprietary models. For example, in our preliminary testing, the model struggled with questions involving triangles and ellipses, potentially pointing to biases in data curation during both pre-training and fine-tuning phases. Furthermore, due to its model size, DeepSeekMath~underperforms relative to GPT-4 in few-shot scenarios—GPT-4 is capable of enhancing its output with few-shot prompts, whereas DeepSeekMath~yields comparable results in zero-shot and few-shot contexts. Looking forward, we aim to further refine our data engineering process to develop an even higher-quality pre-training dataset. We also plan to investigate the future research directions outlined in Section \ref{sec:effec_rl} for optimizing reinforcement learning in large language models.

\end{document}